\documentclass{article}
\usepackage{macros}
\title{Notas de aula - MAT01124}
\author{Leonardo Suarez}
\date{\today}

\begin{document}
\maketitle

Notas de aula da disciplina MAT01124 - Cálculo Numérico I, ministrada pelo prof. Paulo Zingano no semestre letivo 25/2. Aulas ocorrem às terças e quintas-feiras, das 13:30 às 15:10, no laboratório F105. A bibliografia usada pelo prof. é um \emph{mistério}, mas estas notas de aula seguem os primeiro sete capítulos do livro \emph{Cálculo Numérico - Um Livro Colaborativo (Python)}, elaborado pelo REAMAT/UFRGS.

Este curso serve como introdução ao cálculo numérico. Tem como objetivo discutir e aplicar técnicas para resolver problemas matemáticos no computador, com foco na implementação de algoritmos e análise de erros. Nas primeiras três semanas, estudamos aritmética de ponto flutuante, tipos de erros e condicionamento de problemas. Nas semanas 4 a 6, começamos a resolver numericamente equações não lineares com métodos de enquadramento e de ponto fixo. Nas quatro semanas seguintes, vemos problemas de álgebra linear numérica, com foco na solução de sistemas lineares através de métodos diretos e iterativos. Também discutimos métodos para resolução de sistemas não lineares. Por fim, nas semanas 14 a 17, abordamos interpolação e ajuste de curvas.

A avaliação se dará por meio da entrega de dois trabalhos e pela presença. A nota final será dada pela média aritmética das duas maiores notas.

\tableofcontents

\section{Representação de números e aritmética de máquina}

\subsection{Sistema de numeração e mudança de base}

\begin{defn} (sistema de numeração de base \(b\)). 
    Dado um número natural \(b > 1\) e o conjunto de símbolos \(\{\pm, 0, 1, 2, \ldots, b-1\}\), a sequência de símbolos
    \begin{equation}
        (d_n d_{n-1} \ldots d_1 d_0 . d_{-1} d_{-2} \ldots)_b
    \end{equation}
    representa o número positivo

    \begin{equation}
        d_n b^n + d_{n-1} b^{n-1} + \ldots + d_1 b^1 + d_0 b^0 + d_{-1} b^{-1} + d_{-2} b^{-2} + \ldots
    \end{equation}

    Para representar números negativos, utiliza-se o símbolo \(-\) antes da sequência. Cada \(d_i\) é chamado de dígito e deve pertencer ao conjunto \(\{0, 1, 2, \ldots, b-1\}\).
\end{defn}

\begin{rem}
    Se \(b > 10\), é comum utilizar letras do alfabeto para representar dígitos maiores que 9. Por exemplo, na base hexadecimal (base 16), os dígitos são representados como \(0, 1, 2, \ldots, 9, A, B, C, D, E, F\), onde \(A\) representa 10, \(B\) representa 11, e assim por diante até \(F\) que representa 15.
\end{rem}

\subsection{Notação científica e notação normalizada}

\begin{defn}
    Um número \(x\) na base \(b\) é dito estar representado em \emph{notação científica normalizada} quando está escrito na forma

    \begin{equation}
        x = (-1)^s (M)_b \times b^E
    \end{equation}
    onde \((M)_b = (d_0. d_{-1}, d_{-2}, \ldots)_b\) é o mantissa, \(s \in \{0, 1\}\) é o sinal, e \(E\) é o expoente. Além disso, para que a notação esteja normalizada, o dígito mais significativo \(d_0\) deve ser diferente de zero (ou seja, \(d_0 \neq 0\)).
\end{defn}

\subsection{Representação decimal finita}

\begin{defn} (número de dígitos significativos).
    Um número decimal \(x = \pm d_0. d_{-1}\ldots d_{-i}\linebreak d_{-i-1}\ldots d_{-i-n}d_{-i-n-1} \times 10^E\) é dito ter \(n\) \emph{dígitos significativos} quando \(d_j = 0\) para \(j \geq -i\) ou \(j \leq -i-n-1\).
\end{defn}

\subsection{Representação de números em máquina}

\begin{defn}
    A \emph{precisão} \(p\) de uma máquina é o número de dígitos significativos usado para representar um número em sua aritmética.
\end{defn}

\begin{rem}
    Note que \(p = \vert M \vert + 1\) em binário (pois o \(1\) é implícito, logo não é armazenado) e \(p =\vert M \vert\) para outras bases, onde \(\vert M \vert\) é o número de dígitos na mantissa.
\end{rem}

\begin{defn}
    O \emph{épsilon de máquina} \(\epsilon_{mach}\) é definido de forma que \(1+\epsilon_{mach}\) seja o menor número representável maior que 1, isto é, \(1+\epsilon_{mach}\) é representável, mas não existem números representáveis em \((1, 1+\epsilon_{mach})\).
\end{defn}

\subsection{Tipos de erros}

\begin{defn} (erro absoluto e erro relativo).
    Seja \(x\) um número real e \(\overline{x}\) sua aproximação. O \emph{erro absoluto} \(EA\) da aproximação é definido como
    \begin{equation}
        EA := \vert x - \overline{x} \vert
    \end{equation}
    O \emph{erro relativo} \(ER\) é definido como
    \begin{equation}
        ER := \frac{EA}{\vert x \vert}, \quad x \neq 0
    \end{equation}
\end{defn}

\begin{defn} (número de dígitos significativos corretos).
    A aproximação \(\overline{x}\) de \(x\) é dita ter \(k\) \emph{dígitos significativos corretos} se
    \begin{equation}
        ER < \frac{1}{2} \times 10^{-k}
    \end{equation}
\end{defn}

\begin{rem}
    Uma maneira de estimar \(k\) é utilizando a fórmula
    \begin{equation}
        DIGSE(x, \overline{x}) = k = \left\lfloor -\log_{10}(ER) + \frac{1}{2} \right\rfloor
    \end{equation}
\end{rem}

\subsection{Condicionamento de um problema}

\begin{defn}
    Seja \(f\) uma função diferenciável. O \emph{número de condicionamento} de um problema é definido como
    \begin{equation}
        \kappa_f(x) := \left\vert \frac{xf'(x)}{f(x)} \right\vert
    \end{equation}
    \noindent e fornece uma estimativa de quanto os erros na entrada podem afetar a saída do problema.
\end{defn}

\begin{rem}
    De modo geral, quando \(f\) depende de várias variáveis, podemos obter o erro aproximado na saída como
    \begin{equation}
        \delta_f = \vert f(x_1, x_2, \ldots, x_n) - f(\overline{x}_1, \overline{x}_2, \ldots, \overline{x}_n) \vert \approx \sum_{i=1}^{n} \left\vert \frac{\partial f}{\partial x_i} (x_1, x_2, \ldots, x_n)\right\vert \delta_{x_i}
    \end{equation}
\end{rem}

\section{Solução de equações de uma variável}

\subsection{Existência e unicidade}
\subsection{Método da bisseção}
\subsection{Iteração de ponto-fixo}
\subsection{Método de Newton-Raphson}
\subsection{Método das secantes}
\subsection{Critérios de parada}

\section{Solução de sistemas lineares}

\subsection{Eliminação gaussiana}
\subsection{Complexidade de algoritmos em álgebra linear}
\subsection{Sistemas triangulares}
\subsection{Fatoração LU}
\subsection{Método da matriz tridiagonal}
\subsection{Condicionamento de sistemas lineares}
\subsection{Métodos iterativos para sistemas lineares}
\subsection{Cálculo de autovalores e autovetores}

\section{Solução de sistemas de equações não lineares}

\subsection{Método de Newton para sistemas}
\subsection{Linearização de uma função de várias variáveis}

\section{Interpolação}

Dada uma sequência de números reais \(x_1 < \cdots < x_n\), um conjunto de pontos \(\{(x_i, y_i) \in I \times \mathbb{R}\}_{i=1}^{n}\), onde \(I = [x_1, x_n]\) e uma família de funções \(\mathcal{F}_I = \{\varphi : I \to \mathbb{R}\}\), o problema de interpolação consiste em encontrar uma função \(f \in \mathcal{F}_I\) tal que \[f(x_i) = y_i, \quad i = 1, 2, \ldots, n.\]

Chamamos uma tal \(f\) de \emph{função interpoladora} dos pontos dados, ou que \(f\) interpola os pontos dados. Um problema de interpolação cuja família de funções \(\mathcal{F}_I\) é o conjunto dos polinômios de grau no máximo \(n-1\) é chamado de \emph{problema de interpolação polinomial}.

\subsection{Interpolação polinomial}

\begin{thm} (Weierstrass).
    Seja \(f : [a, b] \to \mathbb{R}\) uma função contínua. Para todo \(\delta > 0\), existe um polinômio \(p\) tal que
    \begin{equation}
        \vert f(x) - p(x) \vert < \delta, \quad \forall x \in [a, b].
    \end{equation}
\end{thm}

Para um problema de interpolação polinomial ser bem posto, é necessário restringir o grau dos polinômios. De fato, dado um conjunto de \(n\) pontos \(\{(x_i, y_i)\}_{i=1}^{n}\) com \(x_i \neq x_j\) para \(i \neq j\), escolhemos a família de polinômios \(\mathcal{F} = \mathbb{P}_{n-1}\)\footnote{Isso porque, para garantir a passagem por \(n\), precisamos ajustar exatamente \(n\) coeficientes, dados por um polinômio de grau \(n-1\).}, onde \[\mathbb{P}_{n-1} \defeq \left\{p: x \mapsto \sum_{k=0}^{n-1} a_k x^k : a_0, a_1, \ldots, a_{n-1} \in \mathbb{R}\right\}\] é o conjunto dos polinômios de grau no máximo \(n-1\). 
Assim, o problema de interpolação polinomial consiste em encontrar um \emph{polinômio interpolador} \(p \in \mathbb{P}_{n-1}\) tal que \(p(x_i) = y_i\) para \(i = 1, 2, \ldots, n\).

\begin{thm}
    Seja \(\{(x_i, y_i)\}_{i=1}^{n}\) um conjunto de \(n\) pontos com \(x_i \neq x_j\) para \(i \neq j\). Então, existe um único polinômio \(p \in \mathbb{P}_{n-1}\) tal que \(p(x_i) = y_i\) para \(i = 1, 2, \ldots, n\).
\end{thm}

\begin{proof}[Ideia da demonstração]
    Nota que encontrar os coeficientes \(a_0, \ldots, a_{n-1}\) do polinômio \(p\) é equivalente a resolver um sistema linear com \(n\) equações e \(n\) incógnitas, onde a matriz dos coeficientes é uma \emph{matriz de Vandermonde}, que é invertível quando os \(x_i\) são distintos\footnote{Lembra que o determinante da matriz de Vandermonde é dado pelo produtório duplo \(\prod_{1 \leq i < j \leq n} (x_j - x_i)\)}. Logo, existe um único polinômio interpolador.
\end{proof}

\begin{rem}
    Dado um conjunto de pontos \(\{(x_i,y_i)\}_{i=1}^n\), a existência e unicidade de um polinômio interpolador para esses pontos não são garantidas se \(\deg p \neq n-1\). Uma reta interpola três pontos sse eles são colineares. Por dois pontos passam infinitas parábolas.
\end{rem}

\subsection{Diferenças divididas de Newton}



\subsection{Polinômios de Lagrange}

Dado um conjunto de \(n+1\) pontos \(\{(x_i, y_i)\}_{i=0}^{n}\) com \(x_i \neq x_j\) para \(i \neq j\), a \emph{base de Lagrange} para polinômios de grau no máximo \(n\) é o conjunto \(\{l_0, l_1, \ldots, l_n\}\) dos \emph{polinômios de Lagrange} definidos por \(l_i(x_j) \defeq \delta_{ij}\).
Assim, o polinômio interpolador \(p\) para os pontos dados pode ser escrito como
\begin{equation}
    p(x) = \sum_{i=0}^{n} y_i l_i(x).
\end{equation}

Usando a forma fatorada dos polinômios de Lagrange, podemos obter uma expressão explícita para cada \(l_i\) como
\begin{equation}
    l_i(x) = \prod_{\substack{j=0 \\ j \neq i}}^{n} \frac{x - x_j}{x_i - x_j}.
\end{equation}


\subsection{Aproximação de funções reais por polinômios interpoladores}
\subsection{Interpolação linear segmentada}
\subsection{Interpolação cúbica segmentada - spline}

\section{Ajuste de curvas}

\subsection{Ajuste de uma reta}
\subsection{Ajuste linear geral}
\subsection{Aproximando problemas não lineares por problemas lineares}

\end{document}